% ===================================================================
%  नकसा नं. १४ — सड़क। --- दुकानदार।
%
%  Traduction, faite en 2026, en bhojpouri d'aujourd'hui, sur l'ido
%  de texto/io. Regles de la colonne : voir 10-nakasa-01.tex. Une
%  seule numerotation. 42 blocs, 97 renvois, seize paires a
%  retourner.
%
%  DANS CETTE LANGUE, LE NOM D'UN METIER EST LE NOM D'UNE CASTE, ET
%  CE TABLEAU-CI OBLIGE A LE DIRE. La rue de Rochelle est pleine de
%  boutiquiers, et dans la plaine gangetique presque chaque metier
%  ancien porte un nom qui est aussi celui d'une jati : लोहार le
%  forgeron, बढ़ई le charpentier, कुम्हार le potier, धोबी le
%  blanchisseur, हजाम le barbier, तेली le pressier d'huile, हलवाई le
%  confiseur, दर्जी le tailleur, सोनार l'orfevre, माली le jardinier,
%  रँगरेज le teinturier. Nommer le metier nomme la caste : les deux
%  ne se separent pas.
%
%  ET LA REGLE DE LA COLONNE NE PEUT PAS S'APPLIQUER ICI, IL FAUT LE
%  DIRE PLUTOT QUE DE FAIRE SEMBLANT. Les refus des tableaux 3, 4, 6,
%  9 et 13 — चइत, राम-राम, पतरा, कबड्डी, रोटी, खटिया, मदारी, पचीसी —
%  etaient tous possibles parce qu'il existait AUTRE CHOSE A ECRIRE.
%  Ici il n'y a rien d'autre : दर्जी est le seul mot pour un tailleur,
%  सोनार le seul pour un orfevre, हलवाई le seul pour un confiseur, et
%  les ecarter laisserait les metiers sans nom. La colonne les ecrit
%  donc, comme elle avait deja ecrit लोहार au tableau 5 et धोबिन au
%  tableau 13, sans l'avoir alors remarque. UN REFUS N'EST POSSIBLE
%  QUE LA OU IL Y A QUELQUE CHOSE D'AUTRE A METTRE ; ailleurs ce
%  n'est plus un refus, c'est un trou. Le releve corrige ici ce qu'il
%  laissait croire : la loi « on modernise la langue, jamais les
%  choses » suppose que la langue ait deux mots, et elle n'en a pas
%  toujours deux.
%
%  Les metiers nouveaux, eux, n'ont pas de caste et le montrent : le
%  photographe (20), l'opticien (38), le dactylographe (33), le
%  mecanicien, le chauffeur portent un compose descriptif ou un mot
%  anglais. LE PARTAGE ENTRE LES DEUX EST EXACTEMENT LE PARTAGE ENTRE
%  CE QUI EXISTAIT AVANT LA MACHINE ET CE QUI EST VENU APRES.
%
%  « आया » LA BONNE D'ENFANTS (98) EST LE CINQUIEME MOT PORTUGAIS DE
%  LA COLONNE. Apres मेज (tableau 1), पादरी (3), मिस्त्री (5) et
%  अस्पताल (11), voici aia. Les cinq figuraient dans la liste des onze
%  de la colonne gujaratie : cinq sur cinq franchissent les mille cinq
%  cents kilometres qui separent Diu du Gange, et le tableau 11 avait
%  deja conclu que la couche n'etait pas cotiere. Elle ne l'est
%  decidement pas.
%
%  Enfin le bain turc (2) est « तुर्की हम्माम », et c'est le mot que
%  la colonne levantine emploie pour sa propre salle de bain. Un
%  etablissement parisien de 1926, un mot arabe, et deux colonnes du
%  releve qui l'ecrivent de la meme facon a trois mille kilometres
%  l'une de l'autre.
% ===================================================================

%%K t14-apar-1 apar
\VUsaut{4.0mm}
\VUtitre{15.0pt}{60}{नकसा नं. १४}
\VUsaut{3.0mm}

%%K t14-tit-1 sub
\VUpk{11.4pt}{40}{सड़क। --- दुकानदार।}
\VUsaut{3.0mm}

%%K t14-01-1 p
१. --- हमार दोस्त योहानेस, जे देहात में रहेला, हमरा घरे तीन दिन बितावे
आइल। \VUgras{ओकरा पहिले} ओकरा बड़का सहर देखे के मउका ना मिलल रहे ;
एहीसे हमनी के आपन पूरा दिन घूमे में बितावत बानी जा, आ हम ओकरा के ऊ सब
समझावत बानी जवन ऊ ना जानेला।

%%K t14-02-1 p
२. --- पहिले हमनी के \VUgras{वेधशाला} \textsuperscript{(1)} देखे गइनी जा,
जवन ओकरा के बहुते नीक लागल। लवटत बेरा ओह \VUgras{सड़क} \textsuperscript{(3)}
से गुजरत, जहाँ \VUgras{तुर्की हम्माम} \textsuperscript{(2)} बा, हमनी के एगो
\VUgras{अरथी} \textsuperscript{(4)} से भेंट भइल। \VUgras{सवयातरा} के आगे
\VUgras{पादरी लोग} चलत रहे, आ ओकनी के पाछे \VUgras{मुरदा-गाड़ी}, जवना
में \VUgras{कफन-पेटी} धइल रहे। बहुते दोस्त लोग \VUgras{सोक में डूबल}
परिवार के साथे रहे।

%%K t14-02-2 p
जुलूस के गुजरला के बाद हमनी के ओह बड़का \VUgras{बीयरघर}
\textsuperscript{(5)} के सामने सड़क पार कइनी जा, जवना के छत पर बहुते
\VUgras{ग्राहक} \VUgras{बीयर} पियत रहे, काँच के \VUgras{छाजन}
\textsuperscript{(6)} के नीचे बइठल।

%%K t14-03-1 p
३. --- योहानेस ओह \VUgras{निसान} के देख के दंग रह गइल जवन \VUgras{सराब
बेचे वाला} \textsuperscript{(7)} के दुकान आ \VUgras{बाजा बेचे वाला}
\textsuperscript{(10)} के दुकान के बीच लागल रहे। ऊ हमरा से ओकर मतलब
पूछलस ; हम जवाब देनी कि ऊ \VUgras{अंगरेजी कोंसुलखाना} \textsuperscript{(8)}
के निसान ह, आ हम ओकरा के ऊ \VUgras{कोंसुल} \textsuperscript{(9)} देखवनी जे
बालकनी पर रहलें।

%%K t14-tit-2 sub
\VUpk{10.2pt}{40}{दुकान के सैर।}
\VUsaut{3.0mm}

%%K t14-04-1 p
४. --- जाहिर बा कि हमनी के \VUgras{बाजा} के बसता के सामने रुकनी जा —
\VUgras{हारमोनियम}, \VUgras{ऑरगन} आ \VUgras{पियानो} ; हमनी के
\VUgras{सुरलिपि} आ पियानो के \VUgras{धुन} के चित्र वाला जिल्द देखनी जा,
आ ओह सोना चढ़ल काठ के \VUgras{वीणा} के निहारनी जा जवन निसान के तौर पर
लागल बा।

%%K t14-05-1 p
५. --- \VUgras{झाड़ू वाला} \textsuperscript{(11)} लोग आ \VUgras{झाड़ू के
गाड़ी} \textsuperscript{(12)} से उठल धूर के बादर हमनी के मजबूर कर देलस कि
हमनी के \VUgras{सामान बेचे वाला} \textsuperscript{(13)} के सुंदर
\VUgras{सामान} के सामने आ \VUgras{टीन के मिस्त्री} \textsuperscript{(14)} के
\VUgras{चूल्हा} आ \VUgras{लालटेन} के बसता के सामने से जल्दी गुजर जाईं।

%%K t14-06-1 p
६. --- ऊपर से, इहाँ हमनी के पटरी छोड़े के पड़ल, काहेकि ऊ गठरी से भरल
रहे, जेकरा के \VUgras{ढोवे के गाड़ी} \textsuperscript{(16)} से उतार के ओह
\VUgras{दुकान} \textsuperscript{(15)} में धरल जात रहे जवन अभिए किराया पर
लेहल गइल रहे।

%%K t14-06-2 p
एहसे हमनी के \VUgras{गाड़ी बनावे वाला} \textsuperscript{(17)} के सुंदर
गाड़ी ना देख पवनी जा, बाकिर एहसे हमनी के रुक के ओह \VUgras{पेटी बनावे
वाला} \textsuperscript{(19)} के देखे से ना रुकनी जा, जे आपन पड़ोसी,
\VUgras{साइकिल आ मोटर बेचे वाला} \textsuperscript{(18)}, खातिर एगो पेटी
बनावत रहे। ओकरा बाद, काहेकि थोड़ा देर हो गइल रहे, हमनी के घरे लवट अइनी
जा।

%%K t14-tit-3 sub
\VUpk{10.2pt}{40}{सहर के सैर।}
\VUsaut{3.0mm}

%%K t14-07-1 p
७. --- अगिले दिने हमार माई योहानेस के सुझावली कि ओकर फोटो खिंचावल जाव,
आ ओकरा बाद ऊ हमरा के कहली कि ओकरा के \VUgras{फोटो खींचे वाला}
\textsuperscript{(20)} लगे ले जाईं। खाना के बाद हमनी के ओह कलाकार के
\VUgras{इस्टूडियो} में गइनी जा, जहाँ हमनी के ठीक-ठाक देर तक अगोरे के
पड़ल। मन बहलावे खातिर हमनी के भीत पर टँगल \VUgras{तसवीर}
\textsuperscript{(22)} देखत रहनी जा।

%%K t14-07-2 p
जब हमनी के नंबर आइल, काम में देर ना लागल, आ हमार दोस्त के
\VUgras{कैमरा} \textsuperscript{(21)} के सामने बइठला के तुरते बाद हमनी के
नीचे उतर गइनी जा।

%%K t14-08-1 p
८. --- पहिला मंजिल पर हमनी के \VUgras{हजाम} \textsuperscript{(23)} के खुलल
केवाड़ से देखनी जा कि ओकर छोकड़ा एगो ग्राहक के बाल काटत बा।

%%K t14-08-2 p
सड़क पर पहुँच के हमनी के \VUgras{तमाखू के दुकान} \textsuperscript{(24)} के
सामने से गुजरनी जा। योहानेस ओह \VUgras{लाल लालटेन} \textsuperscript{(25)} आ
ओह \VUgras{मोटका पाइप} से एतना बहल गइल जवन \VUgras{निसान}
\textsuperscript{(26)} के तौर पर लागल बा, कि ऊ दू गो बुढ़उ साहेब से टकरा
गइल जे \VUgras{नसवार सूँघत} \textsuperscript{(27)} बतियावत रहलें।

%%K t14-tit-4 sub
\VUpk{10.2pt}{40}{हलवाई के दुकान।}
\VUsaut{3.0mm}

%%K t14-09-1 p
९. --- \VUgras{सर्राफ} \textsuperscript{(29)} के \VUgras{सो-केस} में सजल
दुनिया भर के \VUgras{नोट} आ \VUgras{सिक्का} कुछ देर हमनी के धेयान खींच
लेलस, बाकिर काहेकि जलपान के बेरा हो गइल रहे, हमनी के बिना जादे रुकले
\VUgras{हलवाई} \textsuperscript{(31)} के इहाँ चल गइनी जा।

%%K t14-10-1 p
१०. --- दुकान में जवन-जवन \VUgras{मिठाई} रहे — \VUgras{केक, मधु के
केक, बिस्कुट} आ हर तरह के \VUgras{बोनबोन} — ऊ सब देख के हमनी के चुने
में मुसकिल भइल। आखिर में योहानेस \VUgras{मलाई के केक} चुनलस, आ हम खाली
एगो छोट \VUgras{चेरी के टार्ट} लेनी, काहेकि मीठ से \VUgras{हमार दाँत
दुखेला}, आ हमरा \VUgras{दाँत के डाक्टर} \textsuperscript{(30)} लगे बार-बार
जाइल बिलकुल ना भावेला।

%%K t14-tit-5 sub
\VUpk{10.2pt}{40}{मसीन से लिखल।}
\VUsaut{3.0mm}

%%K t14-11-1 p
११. --- आपन दोस्त के \VUgras{टाइपराइटर} देखावे खातिर, जेकर बारे में ऊ
सुनले रहे बाकिर देखले ना रहे, हमनी के ओह \VUgras{दफ्तर} \textsuperscript{(32)} में
गइनी जा जवन हमार \VUgras{चचेरा भाई} \textsuperscript{(34)} के ह।
हमनी के ओकरा के आपन \VUgras{सचिव} \textsuperscript{(35)} के चिट्ठी लिखवावत
पवनी जा, जे \VUgras{आसुलिपिक} बा। चिट्ठी खतम होते एगो \VUgras{टाइपिस्ट}
\textsuperscript{(33)} ओकरा के आपन मसीन पर उतार लेलस। आपन रिस्तेदार के काम
में बाधा ना डाले के मन से हमनी के ओकरा लगे खाली कुछे पल रुकनी जा।

%%K t14-12-1 p
१२. --- हमार चचेरा भाई के दफ्तर के नीचे एगो बड़का \VUgras{घड़ीसाज-सोनार}
\textsuperscript{(37)} रहेला, जे ऊपर से सोना के काम भी करेला। ओकर शानदार
दुकान में बहुते \VUgras{घड़ी, पेंडुलम घड़ी, जेब घड़ी}, \VUgras{जंजीर} आ
कीमती \VUgras{जेवर} बा — जइसे \VUgras{मोती के माला}, सोना के
\VUgras{कंगन}, \VUgras{बाली} आ \VUgras{अँगूठी} जवना में \VUgras{नग} आ
\VUgras{हीरा} जड़ल बा। एगो सो-केस खास तौर पर \VUgras{सोना के सामान}
खातिर रखल बा आ ओहमें चाँदी के \VUgras{बरतन} के बड़का संग्रह बा।

%%K t14-13-1 p
१३. --- सड़क के मोड़ पर मुड़त हम देखनी कि घर के \VUgras{नंबर} लगे लागल
\VUgras{तखती} \textsuperscript{(36)} बदल गइल बा। हम जोर से आपन बात कहे
वाला रहनी कि तबहीं \VUgras{फौजी बाजा} के खुसी भरल आवाज सुनाइल। हमनी के
पलटन की ओर दउड़ पड़नी जा, ई सोचले बिना कि \VUgras{ऊ} \VUgras{टोपी वाला}
\textsuperscript{(39)} आ \VUgras{दस्ताना वाला} \textsuperscript{(40)} के दुकान
के सामने से गुजरे वाला बा, आ कि हमनी के ओकर परेड देखे खातिर
\VUgras{चस्मा वाला} \textsuperscript{(38)} के सो-केस के सामने ठाढ़ रह के
भी ओतने नीक जगहा पर रहतीं जा, जवन \VUgras{नाक के चस्मा, कान के चस्मा} आ
\VUgras{छोट दूरबीन} से भरल बा।

%%K t14-tit-6 sub
\VUpk{10.2pt}{40}{परेड।}
\VUsaut{3.0mm}

%%K t14-14-1 p
१४. --- \VUgras{पलटन} \textsuperscript{(41)} के आगे \VUgras{ढोल-नायक}
\textsuperscript{(42)} चलत रहलें, ओकरा पाछे \VUgras{ढोलकिया}
\textsuperscript{(43)}, \VUgras{बिगुल वाला} \textsuperscript{(44)} आ पूरा
\VUgras{बजनियाँ} लोग। \VUgras{जनरल} \textsuperscript{(45)}, जे चउक पर आपन
\VUgras{अरदली} के साथे रहलें, \VUgras{फुहारा के टोंटी}
\textsuperscript{(49)} लगे रुक गइल रहलें, जवन \VUgras{फुहारा}
\textsuperscript{(48)} के ह, ताकि \VUgras{परेड} देख सकसु।

%%K t14-14-2 p
\VUgras{सदर के अफसर} \textsuperscript{(46)} लोग, \VUgras{कर्नल} आ
\VUgras{लेफ्टिनेंट-कर्नल} ओकरा के \VUgras{तलवार} से सलामी देलें।
\VUgras{मेजर} लोग आपन \VUgras{दस्ता} \textsuperscript{(47)} के आगे रहलें आ
हर \VUgras{कप्तान} आपन \VUgras{कंपनी} के हुकुम देत रहलें, जबकि
\VUgras{लेफ्टिनेंट}, \VUgras{उप-लेफ्टिनेंट} आ \VUgras{सूबेदार} लोग आपन
सिपाहियन के बगल में आपन नियमित जगहा पर रहलें।

%%K t14-15-1 p
१५. --- बहुते राहगीर \VUgras{चौराहा} \textsuperscript{(50)} पर जुट गइल
रहलें ; दुकानदार लोग — \VUgras{छाता बनावे वाला} \textsuperscript{(51)},
\VUgras{रँगरेज} \textsuperscript{(53)}, \VUgras{हथियार बनावे वाला}
\textsuperscript{(54)} — \VUgras{सिपाही} लोग के देखे बाहर निकल अइलें। सब
सवारी — \VUgras{किराया के गाड़ी} \textsuperscript{(52)}, \VUgras{असबाब के
गाड़ी} \textsuperscript{(57)} आ \VUgras{पानी छिड़के के गाड़ी}
\textsuperscript{(56)} भी — पटरी के किनारे लाग गइल।

%%K t14-tit-7 sub
\VUpk{10.2pt}{40}{सड़क के दृस्य।}
\VUsaut{3.0mm}

%%K t14-15-2 p
एगो \VUgras{बेपार के दलाल} \textsuperscript{(55)}, हाथ में आपन
\VUgras{नमूना के पेटी} लेके, जल्दी से सड़क पार कइलस, आ \VUgras{ऊपरी छत}
\textsuperscript{(59)} पर बइठल लोग, जवन \VUgras{ओमनिबस}
\textsuperscript{(58)} के ह, नजारा के मजा लेवे खातिर मुड़ गइल। एगो बेचारा \VUgras{घुमंतू गवइया}
\textsuperscript{(60)} के आपन \VUgras{घुमावदार बाजा} \textsuperscript{(61)}
के साथे आपन \VUgras{गीत} रोके के पड़ल, काहेकि ढोल के सोर ओकर आवाज दबा
देलस।

%%K t14-16-1 p
१६. --- जब पलटन \VUgras{कपड़ा के दुकान} \textsuperscript{(64)} के सामने
पहुँचल, \VUgras{सिलाई करे वाली} \textsuperscript{(63)} आ \VUgras{दर्जी}
\textsuperscript{(62)} लोग आपन \VUgras{सिलाई मसीन} आ आपन काम छोड़ के
खिड़की से देखे लागल। \VUgras{दुकानदार} \textsuperscript{(65)}, जे अभिए आपन
बसता पर नया \VUgras{सूट} \textsuperscript{(66)} लगववले रहलें, ओह
\VUgras{ऊनी थान} \textsuperscript{(67)} आ \VUgras{सूती थान} के भीतर धरवावे
के पड़ल जवन पटरी पर \VUgras{नाली के ढक्कन} \textsuperscript{(68)} लगे रखल
रहे, ई डर से कि भीड़ ओकरा के ना गिरा दे ; इहाँ तक कि बुढ़उ \VUgras{फेरी
वाला} \textsuperscript{(69)} के, जेकरा से एगो दरबानिन \VUgras{मोजा} के
भाव-ताव करत रहली, आपन बेचाई पूरा करे खातिर गलियारा में घुसे के पड़ल।

%%K t14-16-2 p
हमनी के पलटन के साथे छावनी तक गइनी जा, \VUgras{आ}, आपन दिन से बहुते
खुस होके, हमनी के खाना के बेरा घरे लवट अइनी जा।

%%K t14-tit-8 sub
\VUpk{10.2pt}{40}{तरह-तरह के धंधा आ पेसा।}
\VUsaut{3.0mm}

%%K t14-17-1 p
१७. --- अगिले दिने हमार बाबूजी हमनी के सुझवलें कि ओहिजा के अखबार के
छापाखाना देखल जाव। हमनी के जोस से मान लेनी जा।

%%K t14-17-2 p
\VUgras{छापे वाला} ऊपर से \VUgras{किताब बेचे वाला} \textsuperscript{(70)},
\VUgras{प्रकासक}, \VUgras{कागज बेचे वाला} आ \VUgras{जिल्दसाज} भी हउवें।
ऊ पहिला मंजिल पर अखबार के \VUgras{संपादक के कमरा} \textsuperscript{(71)}
बनववले बाड़ें, आ ओकरा साथे \VUgras{खोदाई के कमरा} भी, जहाँ \VUgras{लिथो
वाला} \textsuperscript{(72)} आ \VUgras{ताँबा पर खोदे वाला} \textsuperscript{(73)}
लोग काम करेला, आ आखिर में \VUgras{टाइप जोड़े के कमरा}, जहाँ
\VUgras{छपाई मजूर} \textsuperscript{(74)} लोग बा। असली \VUgras{छापाखाना}
\textsuperscript{(75)}, \VUgras{छापे के मसीन} आ दोसर मसीन भुइयाँ-तल पर बा।

%%K t14-tit-9 sub
\VUpk{10.2pt}{40}{योहानेस बहुते सवाल पूछेला।}
\VUsaut{3.0mm}

%%K t14-18-1 p
१८. --- छापाखाना से निकलत योहानेस बहुते अचरज में एगो छोट काँच के
डिब्बा के सामने रुक गइल जवन एगो खंभा पर लागल रहे, \VUgras{गोटा-किनारी
के दुकान} \textsuperscript{(77)} के सामने, आ पूछलस कि ऊ काहे लागल बा।
हमार बाबूजी ओकरा के समझवलें कि ऊ \VUgras{आग के चेतावनी के डिब्बा}
\textsuperscript{(76)} ह। ऊपर से सहर में सब कुछ हमार दोस्त खातिर अचरज के
बात रहे — सादा \VUgras{पाथर के फुहारा} \textsuperscript{(78)} आ ऊ हलुक
\VUgras{ढोवे के तीन पहिया} \textsuperscript{(80)} से लेके, जेकरा के वरदी
पहिनले एगो \VUgras{छोटका नौकर} \textsuperscript{(79)} चलावत रहे,
\VUgras{डाक के गाड़ी} \textsuperscript{(81)} तक।

%%K t14-19-1 p
१९. --- जब हमार बाबूजी कुछ पल खातिर हमनी के छोड़ के \VUgras{कर वसूले
वाला} \textsuperscript{(83)} से बतियावे गइलें, जे एगो \VUgras{वकील}
\textsuperscript{(82)} आ एगो \VUgras{नोटरी} \textsuperscript{(84)} से बतियावे
रुकल रहलें, हमनी के सड़क के आवाजाही आँख से पीछा करत मन बहलवनी जा।
हमनी के सामने तरह-तरह के लोग गुजरल : एगो \VUgras{हलवाई के छोकड़ा}
\textsuperscript{(85)}, जे टोकरी में शानदार \VUgras{मीट-पाई} लेले रहे,
एगो \VUgras{कपड़ा बेचे वाला} \textsuperscript{(86)} के पाछे आइल ; एगो
\VUgras{लालटेन जरावे वाला} \textsuperscript{(87)} एगो \VUgras{गैस वाला}
\textsuperscript{(88)} से बतियावत रहे ; आ एगो \VUgras{नन}
\textsuperscript{(89)}, जे हाथ में एगो अनाथ लइकी के पकड़ले रहली, आपन मठ
की ओर लवटत रहली।

%%K t14-tit-10 sub
\VUpk{10.2pt}{40}{घरे लवटत बेरा।}
\VUsaut{3.0mm}

%%K t14-20-1 p
२०. --- जवना बेरा हमार बाबूजी हमनी से आके मिललें, योहानेस जोर से हँस
पड़ल, दू गो पूरा कालिख में सनल \VUgras{चिमनी साफ करे वाला}
\textsuperscript{(91)} के मैल मुँह आ अजीब फटल कपड़ा देख के।

%%K t14-21-1 p
२१. --- हम \VUgras{फूल बेचे वाली} \textsuperscript{(92)} से आपन माई खातिर
एगो पौधा खरीदल चाहत रहनी, बाकिर हमार बाबूजी हमरा के बतवलें कि ऊ ढोवे
में बहुते भारी होई आ \VUgras{संतरा} लेहल जादे नीक होई। हमनी के
\VUgras{बेचनिहारिन} \textsuperscript{(94)} लगे गइनी जा, आ जब
\VUgras{गवर्नेस} \textsuperscript{(93)}, जे आपन पाले वाला बच्चा खातिर चुनत
रहली, आपन खरीदारी खतम कइली, हमनी के आपन सामान लेनी जा।

%%K t14-22-1 p
२२. --- घरे लवटत हमनी के कई गो जान-पहिचान वाला से भेंट भइल : हमार
परिवार के एगो पुरनकी सहेली, जे आपन \VUgras{साथी मेहरारू}
\textsuperscript{(95)} के बाँह पकड़ले रहली, पुल आ सड़क के
\VUgras{इंजीनियर} \textsuperscript{(96)}, आ हमार एगो चचेरी बहिन आपन
\VUgras{आया} \textsuperscript{(98)} के साथे।

%%K t14-22-2 p
ओकरा बाद हमनी के आपन माई के \VUgras{टोपी बनावे वाली} \textsuperscript{(97)}
से भेंट भइल आ दू गो \VUgras{मठवासी} \textsuperscript{(99)} से, जे सहर में
परदेसी रहलें आ हमार बाबूजी से इस्टेसन के रस्ता पूछे आइल रहलें।
